\section*{Chapter \ref{ch:nw:private-connectivity-to-crypto-venues} --- Private Connectivity to Crypto Venues}

\begin{solution}{exo:nw:private-connectivity-to-crypto-venues:1}
$0.014 \times 730 = 10.22$~USD a month.
\end{solution}

\begin{solution}{exo:nw:private-connectivity-to-crypto-venues:2}
Peering joins the venue's network to each customer's: the venue must coordinate address spaces and routes with every one of them and trusts their
routing. An endpoint service publishes one load balancer to any number of customers, who each get an interface in their own network; the venue's
network stays closed.
\end{solution}

\begin{solution}{exo:nw:private-connectivity-to-crypto-venues:3}
The \omterm{def:nw:trading-in-a-public-cloud:pg}{placement group} decides where the market maker's instances are launched: in the same high-bisection segment as the venue's machines. An endpoint
decides only the route from wherever the instance already is.
\end{solution}

\begin{solution}{exo:nw:private-connectivity-to-crypto-venues:4}
The endpoint has no interface in your zone: every packet will cross a zone boundary. Find out where the venue's gateways are (their zone
identifier); if they are in az1 or az2, move your instances there; if the venue can add an interface in az4 and its gateways are in az4, ask for it.
\end{solution}

\begin{solution}{exo:nw:private-connectivity-to-crypto-venues:5}
It leaves the private network for the provider's edge, terminates at a \omterm{def:nw:the-european-map:pop}{point of presence}, and is carried back to the origin through more equipment
(and usually a TLS termination and a load balancer); in the model that costs \qty{2}{\milli\second} against tens of microseconds for the private
hops.
\end{solution}

\begin{solution}{exo:nw:private-connectivity-to-crypto-venues:6}
Because in that building the \omterm{def:nw:colocation-products-and-how-they-are-sold:mmr}{cross-connect} is metres of fibre straight to the venue's network, while PrivateLink goes out to the cloud provider's
region and back; the venue's own documentation recommends \omterm{def:nw:colocation-products-and-how-they-are-sold:mmr}{cross-connects} for the lowest and most stable latency.
\end{solution}

\begin{solution}{exo:nw:private-connectivity-to-crypto-venues:7}
About \qty{126}{\micro\second} of endpoint hops: the same-zone endpoint's median then reaches plain peering's 99th percentile,
\qty{396}{\micro\second}.
\end{solution}

\begin{solution}{exo:nw:private-connectivity-to-crypto-venues:8}
A \omterm{def:nw:private-connectivity-to-crypto-venues:endpoint}{private endpoint} is private, not necessarily fast: it adds hops, it may be served from another zone, and a shared \omterm{def:nw:trading-in-a-public-cloud:pg}{placement group} or plain peering,
where the venue offers them, are faster. ``Private'' says who can see the traffic, not how long it takes.
\end{solution}

\begin{solution}{pb:nw:private-connectivity-to-crypto-venues:1}
\textbf{Part I.}
\begin{enumerate}
\item 310 and \qty{436}{\micro\second}.
\item 595 and \qty{797}{\micro\second}.
\item \qty{285}{\micro\second} to the median, \qty{360}{\micro\second} to the 99th percentile.
\item Seven times the \qty{40}{\micro\second} step from peering to a same-zone endpoint.
\end{enumerate}
\textbf{Part II.}
\begin{enumerate}[resume]
\item 0.004~USD through the same-zone endpoint, 0.012 through the other.
\item All of it: $0.4~\mathrm{GB} \times 0.01 \times 2$ directions $= 0.008$~USD.
\item $10.22 + 40 = 50.22$~USD and $10.22 + 120 = 130.22$~USD a month.
\item 10.22~USD in one zone, 30.66 in three.
\end{enumerate}
\textbf{Part III.}
\begin{enumerate}[resume]
\item By comparing its instances' zone identifiers with those of the endpoint's interfaces, and by measuring: a median near \qty{600}{\micro\second}
  where 300 was expected is the sign.
\item Move its instances to a zone where the endpoint has an interface (if the venue's gateways are there), ask the venue for an interface in its
  zone, or use another path the venue offers (peering, a shared \omterm{def:nw:trading-in-a-public-cloud:pg}{placement group}).
\item The venue's own segment of the network: about \qty{150}{\micro\second} of median round trip in the model, the best of all.
\item That moving the endpoint helps only if the venue's gateways are in the firm's zone; otherwise the boundary is crossed anyway, and the instances
  must move.
\end{enumerate}
\textbf{Part IV.}
\begin{enumerate}[resume]
\item \emph{Named result}: served from another zone, the \omterm{def:nw:private-connectivity-to-crypto-venues:endpoint}{private endpoint} adds about \qty{285}{\micro\second} to the median round trip (and
  \qty{360}{\micro\second} to the 99th percentile) and 0.008~USD per million 400-byte messages of cross-zone transfer, on top of 0.004~USD of data
  processing.
\item Documentation: the paths, the zone behaviour, the prices. Measurement: chapter~16's round trips. Assumptions: the endpoint hops, the message
  size, the shared-cluster median.
\item Because it costs cents per million messages and tens of dollars a month: its price is negligible next to what the microseconds are worth.
\item The zone identifiers of its gateways and of its endpoint's interfaces, per environment.
\item After every relaunch, after the venue's maintenance, and continuously by measurement (chapter~17).
\item Both are about where the customer's path runs relative to the venue: in a hall the venue equalises the paths; in a cloud region nobody does, and
  the customer must find the short one.
\item The path would be a \omterm{def:nw:colocation-products-and-how-they-are-sold:mmr}{cross-connect} in the building (chapter~9), and the zone question would disappear.
\item It guarantees that traffic stays on the provider's private network, not that it takes the shortest path or stays in one zone.
\end{enumerate}
\end{solution}

\begin{solution}{iq:nw:private-connectivity-to-crypto-venues:1}
Peering routes between two whole networks; PrivateLink exposes one service (behind a load balancer) through an interface in the consumer's network,
without routing between them. Peering has fewer hops; PrivateLink scales to many consumers and keeps the provider's network closed.
\iqlookfor{Routing against service exposure; hops; who controls what.}
\end{solution}

\begin{solution}{iq:nw:private-connectivity-to-crypto-venues:2}
Compare the zone identifiers (not names) of the instances with those of the endpoint's interfaces and the venue's gateways; check the load balancer's
cross-zone setting; and measure: a cross-zone round trip is visibly longer.
\iqlookfor{Identifiers; load-balancer behaviour; measurement.}
\end{solution}

\begin{solution}{iq:nw:private-connectivity-to-crypto-venues:3}
By latency need against cost and eligibility: a shared \omterm{def:nw:trading-in-a-public-cloud:pg}{placement group} or peering if the strategy races, an endpoint if it needs privacy and
reasonable speed, the public path otherwise. Ask for the zones, the rate limits per tier, the gateways' capacity, and the eligibility terms.
\iqlookfor{Tiers mapped to strategies; zones; rate limits; terms.}
\end{solution}

\begin{solution}{iq:nw:private-connectivity-to-crypto-venues:4}
The instance's stack, the endpoint interface, the load balancer, any cross-zone hop, the venue's gateway. Timestamp packets in hardware at the
instance, compare with the venue's acknowledgement timestamps (with care for its clock, chapter~20), and vary one piece at a time (zone, path).
\iqlookfor{Decomposition; hardware timestamps; controlled experiments.}
\end{solution}

\begin{solution}{iq:nw:private-connectivity-to-crypto-venues:5}
It is fair if it is offered on published, non-discriminatory terms to every participant who meets them, as \omterm{def:nw:colocation-products-and-how-they-are-sold:colo}{colocation} must be (chapter~9); it is not
if it is a private arrangement. The question is the venue's conditions, not the technology.
\iqlookfor{Equal conditions within a service; published terms; the \omterm{def:nw:colocation-products-and-how-they-are-sold:colo}{colocation} analogy.}
\end{solution}

\begin{solution}{iq:nw:private-connectivity-to-crypto-venues:6}
Instances in both zones, each connected to the venue's gateways in its own zone (endpoint interfaces or peering in each), with order routing that
prefers the zone of the venue's primary matching engine and fails over to the other; measure both paths continuously and test the fail-over.
\iqlookfor{Zonal alignment; active paths in both zones; fail-over; measurement.}
\end{solution}
